\begin{tikzpicture}[
  x=1cm,
  y=1cm,
  font=\sffamily,
  obs/.style={draw, rounded corners=1pt, fill=gray!10, minimum height=6.5mm,
    align=center, inner xsep=2.5mm, inner ysep=1.4mm, font=\scriptsize},
  tensor/.style={draw, rounded corners=1pt, fill=gray!18, minimum width=61mm,
    minimum height=9mm, align=center, inner sep=1.6mm, font=\scriptsize},
  latent/.style={draw, circle, fill=white, minimum size=8.5mm, align=center,
    inner sep=1pt, font=\small},
  panel/.style={draw=gray!55, rounded corners=2pt, fill=none, inner sep=2.5mm},
  arrow/.style={-{Latex[length=2mm,width=1.3mm]}, line width=.55pt},
  ident/.style={-{Latex[length=2.4mm,width=1.6mm]}, line width=.8pt},
  lab/.style={fill=white, inner sep=1pt, align=center, font=\scriptsize},
  titlelabel/.style={fill=white, inner xsep=1.2mm, inner ysep=.4mm,
    font=\scriptsize\bfseries, anchor=west},
  tinylabel/.style={fill=white, inner sep=.6pt, font=\tiny}
]
  % Stage 1: only the full observed-data law is available.
  \node[obs, minimum width=61mm, minimum height=11mm] (pobs) at (3.6,12.15)
    {$P_{\mathrm{obs}}(O)$\\[-.2mm]
     $O=(X,W,Z,D,Y_D)$};
  \node[titlelabel] at (.35,12.98)
    {識別前：観測データ法則のみ};
  \node[panel, fit=(pobs), minimum width=70mm, minimum height=18mm] (pre) {};
  \node[font=\scriptsize, anchor=south east] at (pre.south east)
    {四角：観測法則};

  % First identification step: recover each supported block law.
  \draw[ident] (pobs.south) -- node[lab]
    {blockwise B-completeness\\命題2.1} (3.6,10.35);

  \node[tensor] (blocklaw) at (3.6,9.75)
    {supported block laws\\[-.2mm]
     $\{\pi_S,\ P(Y_S,Z_S,U_S)\}_{S\in\mathcal S}$};

  \draw[ident] (blocklaw.south) -- node[lab]
    {周辺化・条件付け} (3.6,8.25);

  % The recovered tensors are identified functions of the observed law.
  \node[tensor, minimum height=14mm] (tensor) at (3.6,7.45)
    {回復済みblock tensor\\[-.2mm]
     $T_{abc}(w,x),\quad\{T_{jab}(w,x)\}_{j}$\\[-.2mm]
     $T_{abc}=\sum_f\lambda_f\,M_a\otimes M_b\otimes M_c$};

  % Second identification step: decompose the tensors.
  \draw[ident] (tensor.south) -- node[lab]
    {Kruskal rank $+$ 逐次線形反転\\anchor orderingでlabel固定\quad 定理3.1}
    (3.6,5.72);

  % Stage 2: identified latent measurement structure.
  \node[obs, minimum width=10mm] (w) at (.75,3.75) {$W$};
  \node[obs, minimum width=10mm] (x) at (.75,2.35) {$X$};
  \node[latent] (f) at (3.20,3.75) {$F$};
  \node[obs, minimum width=24mm, minimum height=9mm] (ys) at (6.00,3.75)
    {$Y_a,Y_b,Y_c,Y_j$};

  \draw[arrow] (w) -- node[lab, above] {$Q_F$} (f);
  \draw[arrow] (x) -- (f);
  \draw[arrow] (f) -- node[lab, above] {$Q_j$} (ys);
  \draw[arrow] (x.east) -- ++(.25,0) -| (ys.south);
  \node[font=\scriptsize, align=center] at (5.15,2.80)
    {$Y_a,\ldots,Y_j\ \text{indep.}\mid(F,X)$};

  \node[titlelabel] at (.35,5.42)
    {識別後：tensorと潜在測定構造};
  \node[panel, fit=(blocklaw)(tensor)(w)(x)(f)(ys),
    minimum width=70mm, inner ysep=3.5mm] (post) {};
  \node[font=\scriptsize, anchor=south east] at (post.south east)
    {円：潜在\quad 四角：顕在・識別済み};
\end{tikzpicture}
